% ============================================================
% 一、问题重述
% ============================================================
\section{问题重述}

\subsection{问题背景}
干燥是中药材加工中的重要工序。热风干燥因设备简单、操作方便而应用广泛，但其干燥效率和成品品质易受工艺条件影响\cite{LiWang2024MedicinalDrying,Li2021GastrodiaDrying}。

热风干燥过程中，药材与热空气持续进行热量和水分交换。烘干工艺可分为预热平衡和恒温干燥两个阶段；从干燥动力学看，过程通常表现为升温、恒速和降速三个阶段\cite{Yan2026DryingHeatMassTransfer}。药材内部温度和水分浓度随时间与空间变化，是描述完整烘干过程的重要依据。

随着水分持续流失，药材可能发生收缩并引起几何尺寸变化，进而影响温度和水分分布的预测\cite{MayorSereno2004Shrinkage}。烘房条件、药材参数及尺寸变化共同影响实际烘干过程。

传统依赖反复试验调节工艺参数的方法成本高、周期长。有必要借助数理分析与数值仿真，揭示药材的烘干规律，为过程预测与烘干终点判定提供依据。

\subsection{问题提出}
研究对象为近似圆柱形中药材。根据烘房温湿环境、药材参数及尺寸变化等已知信息，需要解决以下四个问题。

\noindent\textbf{问题一：}建立预热平衡阶段的数学模型，描述药材内部温度和水分浓度随时间、径向位置的变化规律。

\noindent\textbf{问题二：}考虑预热平衡向恒温干燥过渡时的参数变化，将模型拓展至完整烘干过程，并分析初始 $3~\mathrm{h}$ 内温度和水分浓度的演化特征。

\noindent\textbf{问题三：}以药材各处水分浓度均低于 $0.15~\mathrm{kg/kg}$ 作为烘干完成判据，根据水分浓度演化结果确定所需时长。

\noindent\textbf{问题四：}进一步考虑水分流失引起的药材尺寸收缩，并依据半径变化修正模型，重新研究水分迁移过程及烘干时长。
